\scanpage{59}
\mspage{—0.26 bis—}

Putting these two results together, it follows that if $G$ is reductive and the stabilizers are \uline{finite}, $\phi$ being affine is, in essence, a weak topological restriction, i.e. it is equivalent to the properness of $\sigma$. But if the stabilizers are positive dimensional, $\phi$ being affine is a much stronger restriction. At least in characteristic 0, it implies that these stabilizers are themselves reductive groups (cf. Iwahori [\quad]).

The final ``preliminary\raisebox{0.6ex}{''} result is aimed at further motivating the concept of a geometric quotient by showing that, for a free action, it implies all that one might hope for.

\uline{Definition} \uline{0.10}: Let $(Y,\phi)$ be a geometric quotient of $X$ by $G$ (over $S$). Assume that $G$ is flat and of finite type over $S$. Then $X$ is a \uline{principal} \uline{fibre} \uline{bundle} over $Y$, with group $G$, if

\qquad i) $\phi$ is a flat morphism of finite type,

\qquad ii) $\unsure{\bar\Psi|\bar\psi}$ is an isomorphism of $G\mathbin{\underset{S}{x}}X$ and $X\mathbin{\underset{Y}{x}}X$.

\margin{58}
